\section{链式法则就是线性映射复合}
\label{sec:10-Matrix-Calculus/02-Differentials/02}

训练时随手看一眼 profiler，会撞见一个不太合常理的数字：一次反向传播的耗时大约是一次前向的一到两倍。可反向要交出的是几千万个参数各自的偏导数，前向只是把函数值算了一遍。求几千万个导数，凭什么只比算一次函数值贵一倍？

先看老实做法能有多贵。三层复合 \(x\mapsto Ax\mapsto\tanh\mapsto\tfrac12\norm{\cdot-y}_2^2\)，逐分量展开损失对 \(x_i\) 的偏导，写出来是

\[
\frac{\partial\ell}{\partial x_i}=\sum_j\sum_k\frac{\partial\ell}{\partial v_k}\frac{\partial v_k}{\partial u_j}\frac{\partial u_j}{\partial x_i}.
\]

三层就有两重求和，每层几百维，下标已经追不动了，而真实网络有几十层。升一级，改用 Jacobian：链式法则说 \(J_{g\circ f}=J_g J_f\)，逻辑无懈可击。但把尺寸代进去，一个 \(1000\times 1000\) 的线性层，它对输入的 Jacobian 就是 \(A\) 本身，\(10^6\) 个数；后面 \(\tanh\) 的 Jacobian 是对角阵，仍要占一个 \(10^6\) 的位置；再一层又是 \(10^6\)。而你想要的不过是一个与参数同形的矩阵。为什么要途经几个百万级的中间对象才拿到它？

答案在一处很容易被跳过的地方：这串连乘的括号可以随你怎么打。

\begin{center}
\begin{tikzpicture}[x=1cm,y=1cm,font=\footnotesize,>=stealth]
  \node[anchor=west] at (0.0,3.62) {先把 Jacobian 连乘出来，再作用};
  \foreach \x/\lab in {0.2/{\(J_3\)},1.4/{\(J_2\)},2.6/{\(J_1\)}}
    {\draw[black!60,fill=blue!7] (\x,2.4) rectangle (\x+0.9,3.3);
     \node at (\x+0.45,2.85) {\lab};}
  \node at (4.05,2.85) {\(\Longrightarrow\)};
  \draw[black!60,fill=orange!30] (4.85,2.4) rectangle (5.75,3.3);
  \node at (5.30,2.85) {\(d\times d\)};
  \node[anchor=west,black!70] at (6.35,2.85) {每一步乘出一个方阵：\(O(d^3)\)};
  \node[anchor=west] at (0.0,1.62) {带着一列一路乘过去};
  \foreach \x/\lab in {0.2/{\(J_3\)},1.4/{\(J_2\)},2.6/{\(J_1\)}}
    {\draw[black!60,fill=blue!7] (\x,0.4) rectangle (\x+0.9,1.3);
     \node at (\x+0.45,0.85) {\lab};}
  \draw[black!60,fill=blue!16] (3.70,0.4) rectangle (3.95,1.3);
  \node at (3.825,0.12) {\(v\)};
  \node at (4.55,0.85) {\(\Longrightarrow\)};
  \draw[black!60,fill=orange!30] (5.20,0.4) rectangle (5.45,1.3);
  \node[anchor=west,black!70] at (6.35,0.85) {中间结果始终是一列：\(O(d^2)\)};
\end{tikzpicture}

\emph{同一串连乘，挪动括号不改变结果，只改变中间结果的形状；而中间结果的形状就是成本。}
\end{center}

图里两行算的是同一个 \(J_3J_2J_1v\)。上面一行先算 \(J_3J_2\)，中间落下一个 \(d\times d\) 的方阵；下面一行先算 \(J_1v\)，中间落下的始终是一列。结合律保证两个答案完全相同，代价却差着一个维度。反向传播的全部效率来自这一条，与神经网络的结构无关，也与激活函数无关——它只是把括号打在了正确的地方。下面把这笔账算完。

\subsection{导数的复合就是矩阵的连乘}

回到定义。\(f\) 在 \(x\) 处的导数是一个线性映射，把扰动 \(\mathrm dx\) 送成 \(\mathrm dy=Df(x)[\mathrm dx]\)。复合函数的导数就是这两段映射接起来：

\[
D(g\circ f)(x)=Dg(y)\circ Df(x),
\qquad
J_{g\circ f}(x)=J_g(y)\,J_f(x).
\]

先作用的写在右边，与映射的先后次序一致。这条式子不需要单独记忆，一阶近似接一阶近似仍是一阶近似，坐标写法就长这样。

关键在于它还说了一件别的事：映射的复合不要求你先把矩阵造出来。你完全可以让一个向量直接流过这些映射，每一步只做``矩阵乘向量''，从头到尾不构造任何一个完整的中间 Jacobian。链式法则给出的是一个乘积表达式，怎么求值是留给你的自由。

\subsection{括号打在哪一端，代价差一个维度}

把账算清楚。设有 \(L\) 层，每层的 Jacobian 都是 \(d\times d\)，需要的是复合的导数 \(J=J_L\cdots J_1\)。

老实地把 \(J\) 乘出来，要做 \(L-1\) 次矩阵乘矩阵，每次 \(O(d^3)\)，合计 \(O(Ld^3)\)，还得腾出 \(O(d^2)\) 的空间放结果。

换一种打括号的方式。若最终只需要 \(J\) 作用在某个输入方向 \(v\in\R^d\) 上，按 \(J_L(\cdots(J_2(J_1v)))\) 求值，每一步都是矩阵乘向量，\(O(d^2)\)，合计 \(O(Ld^2)\)。这是前向模式，传播的是 JVP。

若需要的是某个输出侧的线性评价 \(u\) 通过 \(J\) 拉回来的结果，按 \(((u^{\trans}J_L)J_{L-1})\cdots J_1\) 求值，每一步是行向量乘矩阵，同样 \(O(d^2)\)，合计 \(O(Ld^2)\)。这是反向模式，传播的是 VJP。

三种做法的结果一模一样，代价相差 \(d\) 倍。深度学习里 \(d\) 动辄成千上万，这一个因子就是能不能训练的分界线。

\subsection{标量损失让反向模式只需要一次}

上面的比较还没说完，因为 JVP 和 VJP 各自只覆盖一个方向。要拿到完整的 Jacobian，前向模式得对 \(n\) 个输入基方向各跑一次，反向模式得对 \(m\) 个输出基方向各跑一次。两者的总代价分别是 \(O(nLd^2)\) 与 \(O(mLd^2)\)，谁划算取决于 \(n\) 和 \(m\) 谁小。

训练神经网络的场景把这个比较推到了极端：\(n\) 是参数个数，可以上亿；\(m\) 是 \(1\)，因为损失是标量。反向模式只需要播下一颗种子 \(\nabla_z\ell\)，一趟走完就拿齐所有参数的梯度，总代价 \(O(Ld^2)\)，与前向计算同阶。开头 profiler 上那个``只贵一倍''的数字就是这么来的：反向做的乘法次数与前向相当，只是每一步换成了转置作用。

代价当然是有的，只是记在另一本账上。反向传播的每一步都要用到前向留下的中间量——\(\tanh\) 的导数要用它的输出，线性层的权重梯度要用它的输入——所以前向必须把激活值存下来，显存占用随深度线性增长。梯度检查点（gradient checkpointing）做的就是拿时间换这块空间：只存一部分激活，反向时把丢掉的那段重算一遍。时间与空间的这笔交换，是选定反向模式之后才出现的问题。

\subsection{输入维数小于输出维数时前向模式更划算}

反向模式的统治地位只在 \(n\gg m\) 时成立，反过来的场景是真实存在的。

一个物理仿真里有五个可调参数，输出是几万个网格点上的场量，你想知道每个输出对这五个参数的灵敏度。此时 \(n=5\)、\(m\) 上万，前向模式跑五趟就够，反向模式要跑上万趟。超参数扫描、少量设计变量的敏感度分析、把一个标量沿轨迹推出去的场合，都属于这一类。

还有一类混合用法值得记住。Hessian-向量积可以写成``先对损失做一次反向拿到梯度，再对这个梯度做一次前向''，前向与反向各一趟，代价与一次普通反向传播同阶，全程不构造 Hessian。下一节会专门算这件事。

判断规则很短：数一数你要的导数有几个输入方向、几个输出方向，选那个方向数少的一端播种。这条规则不涉及任何深度学习的知识，只涉及矩阵连乘的形状。

\subsection{一个两层例子把反向走完}

具体化。取

\[
u=Ax,\qquad v=\tanh u,\qquad \ell=\tfrac12\norm{v-y}_2^2,
\]

\(A\in\R^{5\times 3}\)，\(x\in\R^3\)，\(y\in\R^5\)。前向算出 \(u,v\) 并保留。

反向从损失那一端开始，\(\bar v=\nabla_v\ell=v-y\in\R^5\)。\(\tanh\) 的 Jacobian 是对角阵 \(\diag(1-v_1^2,\dots,1-v_5^2)\)，但没必要把这个 \(5\times 5\) 造出来，对角阵乘向量退化成逐元素乘：

\[
\bar u=\bar v\odot(1-v\odot v)\in\R^5,
\]

代价 \(O(m)\) 而不是 \(O(m^2)\)。逐元素激活函数在反向里之所以近乎免费，原因就在这里。

线性层交出两个梯度。对输入，\(\nabla_x\ell=A^{\trans}\bar u\in\R^3\)。对权重，用上一节的读取程序：\(\mathrm du=(\mathrm dA)x\)，于是

\[
\mathrm d\ell=\bar u^{\trans}(\mathrm dA)x=\trace\bigl(x\bar u^{\trans}\mathrm dA\bigr)
=\trace\bigl((\bar ux^{\trans})^{\trans}\mathrm dA\bigr),
\qquad
\nabla_A\ell=\bar ux^{\trans}\in\R^{5\times 3}.
\]

形状核对通过：\(\bar u\) 是 \(5\times 1\)，\(x^{\trans}\) 是 \(1\times 3\)，外积 \(5\times 3\) 与 \(A\) 同形。要留意形状对只是必要条件，\(\bar ux^{\trans}\) 与 \(x\bar u^{\trans}\) 在 \(5=3\) 时形状同样通过，真正定住答案的是上面那一行迹的整理。整条反向的维数流是 \(\bar v\in\R^5\)、\(\bar u\in\R^5\)、\(\nabla_x\ell\in\R^3\)、\(\nabla_A\ell\in\R^{5\times 3}\)，每一步的形状都由上游维度和当前算子决定，中途不需要任何 reshape。

\subsection{分支要加，共享要累}

计算图不总是一条链。若中间量 \(x\) 同时流向两条路径，\(y=f(x)\)、\(z=g(x)\)、\(\ell=h(y,z)\)，则

\[
\mathrm d\ell=\Bigl(\frac{\partial\ell}{\partial y}Df(x)+\frac{\partial\ell}{\partial z}Dg(x)\Bigr)[\mathrm dx],
\]

\(x\) 的伴随是两路伴随之和。手写计算图时忘记累加共享变量的梯度，是最高频也最难查的错误：每条分支单独看都对，合起来少了一项，训练还照跑，只是收敛到别处去了。残差连接、权重共享、多任务损失都会制造这种分支。

广播和求和是同一条规则的两副面孔。前向把一个值复制到多个位置（例如 bias 加到每一行），反向就要把这些位置的梯度沿广播轴求和；前向把多个元素合成一个（例如损失沿 batch 轴取平均），反向就要把上游那个标量按同样的系数散回每个元素。前向的一对多，对应反向的多对一。

\subsection{链式法则不出声的地方}

有几处地方框架会照常返回一个数，但那个数不是导数。\(\relu\) 在零点不可微，\(\max\)、排序、取整同样如此，自动微分在这些点返回的是一个约定值——右导数、次梯度里的某个特定元素、或某侧极限。凸情形下这个约定值恰好是一个合法次梯度，非凸网络里它就只是一个约定，见\hyperref[sec:09-Convex-Optimization/06-Optimization-Algorithms/03]{``不可微时次梯度为何仍能前进''}。

更隐蔽的偏离来自计算图本身。含随机采样的模型每次前向对应的数学函数都不同，反向拿到的是这一次采样那个函数的梯度，不是期望的梯度。BatchNorm 在训练与推理时走的是两条不同的路径，``执行的程序''与``想求导的函数''之间的对应关系被改动了。\texttt{stop\_gradient} 更直白，它命令自动微分在此处截断链式法则，之后得到的东西对优化可能有用，但它已经不是原复合函数的导数。

所以框架给的``梯度''分两层：链式组合关系是否实现正确，也就是各步 Jacobian 的转置与乘法顺序对不对，这一层框架负责；每个基本算子的局部求导规则是否符合你在这个问题里的语义期望，这一层只能你自己判断。梯度数值检查能抓住第一层的错误，抓不住第二层。

到这里，一阶信息的传播已经算清楚了：括号打对，代价与前向同阶。二阶信息呢？Hessian 是一个 \(d\times d\) 的对象，参数上亿时它连存都存不下。所幸下一节会看到，需要的往往不是这个矩阵，而是它乘一个向量的结果，而这件事同样只要打对括号。
